\documentclass[11pt]{amsart}
\usepackage[margin=1.05in]{geometry}
\usepackage{amsmath,amssymb,amsthm,mathtools}
\usepackage{booktabs,array}
\usepackage[hidelinks]{hyperref}
\hypersetup{pdftitle={Restriction classes and Baker's question on L(1,chi)},pdfauthor={Keio}}

\newtheorem{theorem}{Theorem}[section]
\newtheorem{lemma}[theorem]{Lemma}
\newtheorem{proposition}[theorem]{Proposition}
\newtheorem{corollary}[theorem]{Corollary}
\newtheorem*{mainA}{Theorem A}
\newtheorem*{mainB}{Corollary B}
\newtheorem*{mainC}{Theorem C}
\newtheorem*{mainD}{Theorem D}
\newtheorem*{question}{Baker's question}
\theoremstyle{definition}
\newtheorem{definition}[theorem]{Definition}
\newtheorem{example}[theorem]{Example}
\theoremstyle{remark}
\newtheorem{remark}[theorem]{Remark}

\providecommand{\C}{\mathbb C}
\providecommand{\Q}{\mathbb Q}
\providecommand{\Z}{\mathbb Z}
\providecommand{\F}{\mathbb F}
\newcommand{\Qb}{\overline{\mathbb Q}}
\newcommand{\fq}{\mathfrak q}
\newcommand{\fP}{\mathfrak P}
\newcommand{\Ok}{\mathcal O}
\DeclareMathOperator{\Gal}{Gal}
\DeclareMathOperator{\ord}{ord}
\DeclareMathOperator{\rk}{rank}
\numberwithin{equation}{section}

\begin{document}
\title[Restriction classes and Baker's question]{Restriction classes and Baker's question\\ on the linear independence of $L(1,\chi)$}
\author{Keio}
\date{24 October 2026}

\begin{abstract}
Baker asked in 1975 whether, for every modulus $q$, the values $L(1,\chi)$ at the non-principal Dirichlet characters $\chi$ modulo $q$ are linearly independent over $\Q$.
The answer is known to be affirmative when $\gcd(q,\varphi(q))=1$ (Baker--Birch--Wirsing), and by results of Murty--Murty and Kandhil--Rath the question is equivalent to the $\Q$-linear independence of the algebraic numbers $L(1,\chi)/\pi$ for the odd characters $\chi$.

Let $\lambda(q)$ be the exponent of $(\Z/q\Z)^\times$, $\delta=\gcd(q,\lambda(q))$, and $H$ the kernel of $(\Z/q\Z)^\times\to(\Z/\delta\Z)^\times$.
We show that every $\Q$-linear relation among the odd values $L(1,\chi)$ splits along the fibres of the restriction map $\chi\mapsto\chi|_H$, and that inside a fibre the values are proportional over $\Q(\zeta_{\lambda(q)})$; each fibre has $\max(1,\varphi(\delta)/2)$ elements.
Consequently Baker's question has an affirmative answer whenever $\varphi(\delta)\le2$, which extends the known criteria to infinitely many further moduli, and in general the odd values span a space of dimension at least $\varphi(q)/\max(2,\varphi(\delta))$.

For prime squares $q=p^2$ we write the fibres as Fermat-quotient twisted cotangent sums, that is, as values of Iwasawa power series at $p$-th roots of unity, decompose them into Galois-isotypic pieces, and separate the pieces $p$-adically.
This answers Baker's question for $q=p^2$ whenever $p-1$ is $2^k$, $2\ell$ or $4\ell$ with $\ell$ an odd prime; the last two cases use the class number of $\Q(\sqrt{-p})$ and the parity of the relative class number of the imaginary cyclic quartic field of conductor $p$.
Under an explicit condition on Bernoulli numbers and Fermat quotients modulo $p$, which we verify for all $p<1.6\cdot10^5$, the same holds whenever $p-1=2^v\ell$.
We also give a rational criterion for prime squares and certify an affirmative answer for every $q\le10^4$ and for $q=p^2$ with $p<5000$.

The question remains open in general.
We explain why the remaining fibres of prime-square moduli are controlled by higher coefficients of Iwasawa power series modulo $p$, which known results do not control.
\end{abstract}
\maketitle
\tableofcontents
\section{Introduction}

For a modulus $q\ge3$ let $X_q$ be the set of non-principal Dirichlet characters modulo $q$.
For a prime $p$, Chowla \cite{Cho70} proved that the numbers $\cot(\pi a/p)$, $1\le a\le (p-1)/2$, are linearly independent over $\Q$.
Answering a question of Chowla, Baker, Birch and Wirsing \cite{BBW73} showed that the numbers $L(1,\chi)$, $\chi\in X_q$, are linearly independent over $\Q$ whenever $\gcd(q,\varphi(q))=1$, in particular for every prime $q$.
In his book \cite[p.~48]{Bak75} Baker remarked that it would be of interest to know whether this holds for every modulus.

\begin{question}
Are the numbers $L(1,\chi)$, $\chi\in X_q$, linearly independent over $\Q$ for every $q\ge3$?
\end{question}

The question splits according to parity.
For even $\chi$ the value $L(1,\chi)$ is a linear form in logarithms of cyclotomic units with algebraic coefficients, and Murty and Murty \cite{MM11} proved, using Baker's theorem and Ramachandra's units, that the even values are even linearly independent over $\Qb$.
For odd $\chi$ the value $L(1,\chi)$ is an algebraic multiple of $\pi$, so the odd values span a one-dimensional $\Qb$-space; Kandhil and Rath \cite{KR22} showed, again by Baker's theorem, that the $\Q$-spans of the even and of the odd values intersect trivially.
Hence Baker's question is equivalent to the $\Q$-linear independence of the \emph{algebraic} numbers $L(1,\chi)/\pi$, $\chi$ odd; this part is open in general \cite{KR22}, and since these numbers are algebraic, transcendence methods give no direct information about it.
Besides \cite{BBW73}, the known results rest on Okada's theorem \cite{Oka80} on the linear independence of cotangents over number fields $F$ with $F\cap\Q(\zeta_q)=\Q$; the resulting criterion, implicit in the work of Gun and Kandhil \cite[Theorem~2]{GK22} on unions over pairwise coprime moduli, answers Baker's question when $\gcd(q,\varphi(q))\le2$.

In this paper we isolate the arithmetic structure of the possible relations, prove new cases of Baker's question, among them infinitely many moduli of a simple shape, and verify it in a large range.
We write $U_q=(\Z/q\Z)^\times$, $\lambda=\lambda(q)$ for the exponent of $U_q$ (Carmichael's function), and
\[
\delta=\gcd(q,\lambda(q)),\qquad H=\ker\bigl(U_q\to U_\delta\bigr)=\{a\in U_q:\ a\equiv1\pmod\delta\},\qquad K=\Q(\zeta_\lambda).
\]
For a character $\mu$ of $H$ let $C_\mu$ be the set of odd characters $\chi$ modulo $q$ with $\chi|_H=\mu$; we call these sets the \emph{restriction classes}.

\begin{mainA}
Let $q\ge3$.
\begin{enumerate}
\item If $x_\chi\in K$ ($\chi$ odd modulo $q$) satisfy $\sum_\chi x_\chi L(1,\chi)=0$, then
\[
\sum_{\chi\in C_\mu}x_\chi L(1,\chi)=0\qquad\text{for every character $\mu$ of $H$.}
\]
\item If $\chi,\chi'\in C_\mu$, then $L(1,\chi')/L(1,\chi)\in K$.
\item Every non-empty $C_\mu$ is a coset $\chi E_\delta$ of the group $E_\delta$ of even characters of $U_q$ that are trivial on $H$, which is the group of even characters of $U_\delta$; it has one element if $\delta\le2$ and $\varphi(\delta)/2$ elements if $\delta\ge3$.
\end{enumerate}
Consequently Baker's question for $q$ is equivalent to the following statement: for every restriction class $C_\mu$ and one (any) $\chi_\mu\in C_\mu$, the numbers $L(1,\chi)/L(1,\chi_\mu)\in K$, $\chi\in C_\mu$, are linearly independent over $\Q$.
\end{mainA}

\begin{mainB}
If $\varphi(\delta)\le2$, that is $\delta\in\{1,2,3,4,6\}$, then Baker's question has an affirmative answer for $q$; in this case the odd values are even linearly independent over $\Q(\zeta_{\lambda})$.
For every $q\ge3$ with $\delta\ge3$ the odd values $L(1,\chi)$ span a $\Q$-space of dimension at least $\varphi(q)/\varphi(\delta)$.
\end{mainB}

The condition $\varphi(\delta)\le2$ is much weaker than $\gcd(q,\varphi(q))=1$: it holds for instance for $q=8,9,12,16,18,24,36$, for all $q=2p,\,4p$ and $3p$ with $p$ prime, for $q=8p$ with $p\not\equiv1\pmod 8$, and for $q=9p$ with $p\ne3$ and $p\not\equiv1\pmod 9$.
Corollary B is elementary: it also has a short direct proof along the lines of Okada's argument (Remark~\ref{rem:direct}), and the case $\delta\le2$ is Okada's theorem applied over $\Q(\zeta_\lambda)$; we have not found the cases $\delta\in\{3,4,6\}$ stated in the literature.
Among the $9998$ moduli $3\le q\le 10^4$ it holds for $7290$, whereas $\gcd(q,\varphi(q))\le2$ holds for $4683$ of them (see Remark~\ref{rem:compare}).

Corollary B leaves the moduli with $\varphi(\delta)\ge4$, where the restriction classes have at least two elements and the question becomes a genuinely new one about a few numbers in $\Q(\zeta_\lambda)$.
The first such moduli are $q=25,27,32,49,50,54,55$.
For prime squares we make the classes explicit.
Let $p$ be an odd prime and $\fq(a)=(a^{p-1}-1)/p\bmod p$ the Fermat quotient.
For $c\in\F_p^\times$ the characters $a\mapsto\zeta_p^{c\fq(a)}$ of $U_{p^2}$ have order $p$, and the non-trivial restriction classes modulo $p^2$ are the sets $\{\theta\cdot\zeta_p^{c\fq}:\ \theta \text{ odd modulo } p\}$.
Up to a common factor, $L(1,\theta\zeta_p^{c\fq})$ is
\[
\theta(c)\sum_{b=1}^{p-1}\bar\theta(b)\,\frac{\zeta_p^{-c\fq(b)}}{\zeta_p^{\,c b^{-1}}-1}\qquad(b^{-1}\text{ the inverse of $b$ modulo $p$}),
\]
a cotangent sum twisted by the Fermat quotient; $p$-adically it is, up to the root of unity $-\theta(c)$, the value at $\zeta_p^{-c}-1$ of the Iwasawa power series of the $p$-adic $L$-function attached to $\bar\theta\omega$, $\omega$ the Teichm\"uller character (Remark~\ref{rem:iwasawa}).
We split the classes into Galois-isotypic pieces and separate the pieces by $\fP$-adic valuations, using a polar term that comes from the Teichm\"uller character together with the first two coefficients of the $\fP$-adic expansions.

\begin{mainC}
Let $p$ be an odd prime. Baker's question has an affirmative answer for $q=p^2$ in each of the following cases:
\begin{enumerate}
\item $p$ is a Fermat prime;
\item $p=2\ell+1$ with $\ell$ an odd prime;
\item $p=4\ell+1$ with $\ell$ an odd prime;
\item $p=2^v\ell+1$ with $v\ge3$ and $\ell$ an odd prime, provided an explicit condition $(\mathrm C_p)$ on Bernoulli numbers and Fermat quotients modulo $p$ holds (Theorem~\ref{thm:psquare}); this condition holds for every such $p<1.6\cdot10^5$.
\end{enumerate}
\end{mainC}

Cases (1)--(3) are unconditional.
Besides the polar term they use the class number $h(-p)$ in case (2), and in case (3) the fact that the relative class number of the imaginary cyclic quartic field of conductor $p$ is odd, which forces both coordinates of $B_{1,\theta}$, $\theta$ quartic, to be odd.
Cases (2) and (3) are expected to contain infinitely many primes, but this is not known.
In case (4) the constant terms alone do not suffice (some of them vanish, for instance, for $p=41$), and $(\mathrm C_p)$ also involves the next coefficient, a discrete derivative of $B_n/n$ in the $p$-adic direction.
For general $p$ the descent produces up to $\tau(m)$ pieces to be separated, where $p-1=2^vm$ with $m$ odd; when $m$ is $1$ or a prime there are at most two, and the two available coefficients separate them for $v\le2$, and for $v\ge3$ whenever $(\mathrm C_p)$ holds; in general the pieces are governed by higher coefficients modulo $p$ of Iwasawa power series, which known results do not control (Remark~\ref{rem:iwasawa}).

\begin{mainD}
Baker's question has an affirmative answer for every $q\le 10^4$, and for $q=p^2$ for every prime $p<5000$.
\end{mainD}

Among the $9998$ moduli $3\le q\le10^4$, the $7290$ with $\varphi(\delta)\le2$ are covered by Corollary B; for the remaining $2708$ moduli, every restriction class ($964{,}083$ classes in all, of sizes up to $729$ for $q=3^8$) is certified by a nonsingular matrix of embeddings into a finite field (Section~\ref{sec:comp}).
For prime squares we use the sufficient condition of Theorem~\ref{thm:star}: the linear independence of $(p-1)/2$ explicit cotangents twisted by Fermat quotients.

We stress what is not proved.
The general question remains open: Theorem A reduces it to classes of $\varphi(\delta)/2$ numbers in $\Q(\zeta_\lambda)$, but we have no general argument for these classes; Theorem C covers only special prime squares, and the families in (2)--(4) are not known to be infinite.
Section~\ref{sec:remarks} describes the two kinds of classes that remain, those attached to primes $p\parallel q$ dividing $\delta$ (``tame'' classes, partly accessible through Stickelberger's theorem) and those attached to $p^2\mid q$ (``wild'' classes).

The paper is organised as follows.
Section~\ref{sec:prelim} collects the standard facts we use.
Theorem A and Corollary B are proved in Section~\ref{sec:classes}.
Section~\ref{sec:psquare} treats prime squares and proves Theorem C.
The computations behind Theorem D are described in Section~\ref{sec:comp}, and Section~\ref{sec:remarks} contains remarks on the remaining cases.
\section{Preliminaries}\label{sec:prelim}

Throughout, $q\ge3$, $U_q=(\Z/q\Z)^\times$, $\lambda=\lambda(q)$ is the exponent of $U_q$, $\delta=\gcd(q,\lambda)$, $N=\operatorname{lcm}(q,\lambda)$, and $\zeta_n=e^{2\pi i/n}$.
Characters modulo $q$ are characters of $U_q$, extended by $0$; their values lie in $\mu_\lambda\subset\Q(\zeta_\lambda)$.
For a primitive character $\chi$ of conductor $f$ we use the Gauss sum $\tau(\chi)=\sum_{x\bmod f}\chi(x)\zeta_f^x$ and the generalised Bernoulli number $B_{1,\chi}=f^{-1}\sum_{a=1}^{f}\chi(a)a$.

\subsection{Odd values as cotangent sums}
Let $f$ be a $q$-periodic function with $f(-n)=-f(n)$.
Then $\sum_{n=1}^qf(n)=0$, the series $L(1,f)=\sum_{n\ge1}f(n)/n$ converges, and
\begin{equation}\label{eq:cot}
L(1,f)=\frac{\pi}{2q}\sum_{a=1}^{q-1}f(a)\cot\frac{\pi a}{q}.
\end{equation}
Indeed $L(1,f)=-q^{-1}\sum_{a=1}^{q}f(a)\psi(a/q)$ with the digamma function $\psi$, and pairing $a$ with $q-a$ and using $\psi(1-x)-\psi(x)=\pi\cot(\pi x)$ gives \eqref{eq:cot}.
For $a\in U_q$ put
\[
y_a=\frac{\zeta_q^a+1}{\zeta_q^a-1}=-i\cot\frac{\pi a}{q}\in\Q(\zeta_q),
\]
and for an odd character $\chi$ modulo $q$
\[
W_\chi=\sum_{a\in U_q}\chi(a)\,y_a\in\Q(\zeta_N),\qquad\text{so that}\qquad L(1,\chi)=\frac{\pi i}{2q}\,W_\chi .
\]
Since $L(1,\chi)\neq0$ for every non-principal $\chi$, we have $W_\chi\neq0$.
Thus Baker's question for the odd characters asks whether the numbers $W_\chi$, $\chi$ odd, are linearly independent over $\Q$.
Note also that $q\,y_a=q+2\sum_{j=0}^{q-1}j\zeta_q^{aj}\in\Z[\zeta_q]$, because $\sum_{j=0}^{q-1}j\omega^j=q/(\omega-1)$ for every $q$-th root of unity $\omega\neq1$; hence $qW_\chi\in\Z[\zeta_N]$.

For later use we recall that for an odd primitive character $\chi$ of conductor $f$
\begin{equation}\label{eq:Lprim}
L(1,\chi)=\frac{\pi i\,\tau(\chi)}{f}\,B_{1,\bar\chi},
\end{equation}
and that for an odd character $\chi$ modulo $q$ induced by $\chi^*$ of conductor $f$ we have $L(1,\chi)=L(1,\chi^*)\prod_{r\mid q,\ r\nmid f}(1-\chi^*(r)/r)$, the product over primes $r$.

\subsection{The Galois action}
Let $u\in U_N$ and let $\sigma_u\in\Gal(\Q(\zeta_N)/\Q)$ be given by $\zeta_N\mapsto\zeta_N^u$.
Write $s=u\bmod q\in U_q$ and $t=u\bmod\lambda\in U_\lambda$.
Since $\sigma_u(\chi(a))=\chi(a)^t$ and $\sigma_u(y_a)=y_{sa}$, substituting $b=sa$ gives
\begin{equation}\label{eq:galois}
\sigma_u(W_\chi)=\overline{\chi^t}(s)\,W_{\chi^t}.
\end{equation}
By the Chinese remainder theorem the pairs $(s,t)$ that occur are exactly those with $s\equiv t\pmod\delta$.

\subsection{Reduction to odd characters}
Murty and Murty \cite{MM11} (see also \cite[Theorem~3]{KR22}) proved that the values $L(1,\chi)$, $\chi$ even and non-principal modulo $q$, are linearly independent over $\Qb$.
Kandhil and Rath \cite[Theorem~5]{KR22} proved that the $\Q$-span of the even values meets the $\Q$-span of the odd values only in $0$.
(For the reader's convenience: an element of the intersection would give $\sum_{\chi\ \mathrm{even}}x_\chi L(1,\chi)=\alpha\pi$ with $\alpha$ algebraic; the left side is a $\Qb$-linear combination of logarithms of positive real algebraic numbers, and $\pi=-i\log(-1)$, so Baker's theorem forces both sides to vanish.)
Consequently:

\begin{lemma}\label{lem:oddpart}
Baker's question has an affirmative answer for $q$ if and only if the numbers $W_\chi$, $\chi$ odd modulo $q$, are linearly independent over $\Q$.
\end{lemma}
\section{Restriction classes}\label{sec:classes}

We keep the notation of Section~\ref{sec:prelim} and put $K=\Q(\zeta_\lambda)$ and
\[
\Gamma=\Gal\bigl(\Q(\zeta_N)/K\bigr)=\{\sigma_u:\ u\in U_N,\ u\equiv1\pmod\lambda\}.
\]

\begin{lemma}\label{lem:Gamma}
The map $\sigma_u\mapsto u\bmod q$ is an isomorphism of $\Gamma$ onto $H=\{s\in U_q:\ s\equiv1\pmod\delta\}$.
For $\sigma_u\in\Gamma$ with $u\equiv s\pmod q$ and every odd $\chi$ we have $\sigma_u(W_\chi)=\bar\chi(s)W_\chi$.
\end{lemma}

\begin{proof}
An element $u\in U_N$ is determined by its residues modulo $q$ and modulo $\lambda$, since $N=\operatorname{lcm}(q,\lambda)$; so the map is injective.
By the Chinese remainder theorem the system $u\equiv s\pmod q$, $u\equiv1\pmod\lambda$ is solvable if and only if $s\equiv1\pmod{\gcd(q,\lambda)}$, which identifies the image with $H$.
The second assertion is \eqref{eq:galois} with $t=1$.
\end{proof}

\begin{lemma}\label{lem:eigen}
For a character $\mu$ of $H$ let $V_\mu=\{x\in\Q(\zeta_N):\ \sigma_u(x)=\bar\mu(u\bmod q)\,x\text{ for all }\sigma_u\in\Gamma\}$.
Then each $V_\mu$ is a $K$-vector space of dimension one, $\Q(\zeta_N)=\bigoplus_\mu V_\mu$, and $W_\chi\in V_{\chi|_H}$ for every odd $\chi$.
The projection $\Q(\zeta_N)\to V_\mu$ is $e_\mu=|H|^{-1}\sum_{\sigma_u\in\Gamma}\mu(u\bmod q)\,\sigma_u$, and it is $K$-linear.
\end{lemma}

\begin{proof}
The exponent of $H\subseteq U_q$ divides $\lambda$, so the characters of $H\cong\Gamma$ take values in $K$, and the $e_\mu$ are orthogonal idempotents in $K[\Gamma]$ with sum $1$; they are $K$-linear because $\Gamma$ fixes $K$.
By the normal basis theorem $\Q(\zeta_N)$ is a free $K[\Gamma]$-module of rank one, so $e_\mu\Q(\zeta_N)=V_\mu$ has $K$-dimension one.
The last assertion is Lemma~\ref{lem:Gamma}.
\end{proof}

\begin{proof}[Proof of Theorem A]
(1) Since $L(1,\chi)=\frac{\pi i}{2q}W_\chi$, a relation $\sum_\chi x_\chi L(1,\chi)=0$ with $x_\chi\in K$ is a relation $\sum_\chi x_\chi W_\chi=0$.
Applying the $K$-linear projection $e_\mu$ and using $e_\mu W_\chi=W_\chi$ if $\chi|_H=\mu$ and $e_\mu W_\chi=0$ otherwise (Lemma~\ref{lem:eigen}) gives $\sum_{\chi\in C_\mu}x_\chi W_\chi=0$.

(2) For $\chi,\chi'\in C_\mu$ the numbers $W_\chi,W_{\chi'}$ are non-zero elements of the one-dimensional $K$-space $V_\mu$, so their quotient lies in $K$; it equals $L(1,\chi')/L(1,\chi)$.

(3) Two characters agree on $H$ if and only if their quotient is trivial on $H$, that is, factors through $U_q/H$, which reduction modulo $\delta$ identifies with $U_\delta$.
If $\chi,\chi'$ are odd, then $\chi'\bar\chi$ is even; conversely $\chi\eta$ is odd and agrees with $\chi$ on $H$ for every even character $\eta$ of $U_\delta$.
Hence $C_\mu=\chi E_\delta$ for any $\chi\in C_\mu$.
If $\delta\le2$ the group $U_\delta$ is trivial; if $\delta\ge3$ then $-1\neq1$ in $U_\delta$ and exactly half of the $\varphi(\delta)$ characters of $U_\delta$ are even.

The final statement follows from (1), (2) and Lemma~\ref{lem:oddpart}: a $\Q$-linear relation among the odd values is a sum of relations inside the classes, and inside $C_\mu$ a relation $\sum x_\chi L(1,\chi)=0$ is the same as $\sum x_\chi\,L(1,\chi)/L(1,\chi_\mu)=0$.
\end{proof}

\begin{proof}[Proof of Corollary B]
If $\varphi(\delta)\le2$, every restriction class has one element by Theorem~A(3), and Theorem~A(1) shows that the odd values are linearly independent over $K$, in particular over $\Q$; Lemma~\ref{lem:oddpart} gives the affirmative answer.
In general the $\Q$-spans of the different classes form a direct sum, by Theorem~A(1), and each class contributes at least one dimension.
The $\varphi(q)/2$ odd characters fall into classes of $\varphi(\delta)/2$ elements when $\delta\ge3$, so there are $\varphi(q)/\varphi(\delta)$ classes.
\end{proof}

\begin{remark}[The class elements as a Fourier transform]\label{rem:fourier}
Fix $\chi_0\in C_\mu$ and put
\[
Y_b=\sum_{a\in U_q,\ a\equiv b\ (\mathrm{mod}\ \delta)}\chi_0(a)\,y_a\qquad(b\in U_\delta).
\]
Then $W_{\chi_0\eta}=\sum_{b\in U_\delta}\eta(b)Y_b$ for every even character $\eta$ of $U_\delta$, and all $Y_b$ lie in the line $V_\mu$.
So a restriction class is the Fourier transform over $U_\delta$ of the $\varphi(\delta)$ vectors $Y_b$ of one $K$-line, and Baker's question for $C_\mu$ asks whether no non-zero even function $h=\sum_\eta x_\eta\eta$ with \emph{rational} coefficients $x_\eta$ satisfies $\sum_bh(b)Y_b=0$.
For $K$-valued coefficients such functions always exist once $\varphi(\delta)\ge4$, which is why rationality has to enter.
\end{remark}

\begin{remark}[Comparison with earlier criteria]\label{rem:compare}
Since $\lambda\mid\varphi(q)$ and both have the same prime factors, $\delta=1$ is equivalent to $\gcd(q,\varphi(q))=1$, the hypothesis of \cite{BBW73}.
Okada's theorem \cite{Oka80} asserts that the numbers $\cot(\pi a/q)$, $1\le a<q/2$, $\gcd(a,q)=1$, are linearly independent over every number field $F$ with $F\cap\Q(\zeta_q)=\Q$; as the values $\chi(a)$ lie in $\Q(\zeta_{\varphi(q)})$, it yields Baker's question under the hypothesis $\Q(\zeta_{\varphi(q)})\cap\Q(\zeta_q)=\Q$, that is $\gcd(q,\varphi(q))\le2$; this is implicit in \cite[Theorem~2]{GK22}.
Applied with $F=\Q(\zeta_\lambda)$, for which $F\cap\Q(\zeta_q)=\Q(\zeta_\delta)$, it yields the case $\delta\le2$ of Corollary B.
We have not found the cases $\delta\in\{3,4,6\}$ in the literature; they occur for instance for $q=9,16,18,36,48$, for $q=9p$ with $p\ge5$ prime and $p\not\equiv1\pmod 9$ ($\delta=3$), and for $q=16p$ with $p$ an odd prime, $p\not\equiv1\pmod8$ ($\delta=4$).
In explicit terms, $\delta\in\{1,2,3,4,6\}$ holds if and only if: every prime $p\ge5$ dividing $q$ divides it exactly once and divides no $r-1$ with $r\mid q$ prime; if $9\mid q$ then $27\nmid q$ and no prime $r\mid q$ satisfies $r\equiv1\pmod9$; if $8\mid q$ then $32\nmid q$ and no prime $r\mid q$ satisfies $r\equiv1\pmod 8$; and not both $3\mid\delta$ and $4\mid\delta$.
Among the $9998$ moduli $3\le q\le10^4$ there are $3112$ with $\gcd(q,\varphi(q))=1$, $4683$ with $\gcd(q,\varphi(q))\le2$, $5143$ with $\delta\le2$, and $7290$ with $\varphi(\delta)\le2$; thus $2147$ of them have $\delta\in\{3,4,6\}$.
\end{remark}

\begin{remark}[A direct argument for Corollary B]\label{rem:direct}
Corollary B can also be proved without Theorem A, along the lines of Okada's argument.
Let $A$ be a set of representatives of $U_q/\{\pm1\}$.
The $\varphi(q)/2$ numbers $y_a$, $a\in A$, are linearly independent over $\Q$ (a rational relation, transported by $\Gal(\Q(\zeta_q)/\Q)$ and Fourier transformed, would give $\hat c(\bar\chi)W_\chi=0$ for every odd $\chi$), so they form a $\Q$-basis of the $(-1)$-eigenspace $\Q(\zeta_q)^-$ of complex conjugation.
If $\delta\in\{3,4,6\}$, the field $D=\Q(\zeta_\delta)\subseteq\Q(\zeta_q)$ is imaginary quadratic, so $D\cdot\Q(\zeta_q)^-=\Q(\zeta_q)^-\oplus\sqrt{-d}\,\Q(\zeta_q)^-=\Q(\zeta_q)$, and the $y_a$ are linearly independent over $D$; if $\delta\le2$ put $D=\Q$.
Since $\Q(\zeta_\lambda)\cap\Q(\zeta_q)=\Q(\zeta_\delta)=D$ and both fields are Galois, they are linearly disjoint over $D$, so the $y_a$ remain linearly independent over $\Q(\zeta_\lambda)$, and so do the $W_\chi=2\sum_{a\in A}\chi(a)y_a$, $\chi$ odd.
This argument, however, gives nothing when $\varphi(\delta)\ge4$, whereas Theorem A describes the possible relations for every $q$.
\end{remark}

\subsection{Galois symmetry inside a class}
Let $P=\{t\in U_\lambda:\ t\equiv1\pmod\delta\}$.
For $t\in P$ choose $u_t\in U_N$ with $u_t\equiv1\pmod q$ and $u_t\equiv t\pmod\lambda$; by \eqref{eq:galois}, $\sigma_{u_t}(W_\chi)=W_{\chi^t}$.

\begin{lemma}\label{lem:Psym}
If $\sum_\chi x_\chi W_\chi=0$ with $x_\chi\in\Q$, then $\sum_\chi x_\chi W_{\chi^t}=0$ for every $t\in P$.
The group $P$ permutes the restriction classes, $C_\mu\mapsto C_{\mu^t}$, and the stabiliser $P_\mu$ of $C_\mu$ acts on the $\Q$-space of rational relations inside $C_\mu$.
\end{lemma}

\begin{proof}
Apply $\sigma_{u_t}$, which fixes $\Q$; and $(\chi^t)|_H=(\chi|_H)^t$.
\end{proof}

In Section~\ref{sec:psquare} we use Lemma~\ref{lem:Psym} for $q=p^2$, where $P_\mu=P\cong U_{p-1}$ acts on every class through the tame parts of the characters.
\section{Prime squares}\label{sec:psquare}

Throughout this section $p\ge5$ is a prime and $q=p^2$.
Then $U_{p^2}$ is cyclic of order $p(p-1)$, $\lambda=p(p-1)$, $\delta=p$, $N=p^2(p-1)$, and $H=\{a\equiv1\pmod p\}$ has order $p$.
(For $p=3$ we have $\varphi(\delta)=2$, and $q=9$ is covered by Corollary B.)
Integers $b$ with $1\le b\le p-1$ are used as representatives of $\F_p^\times$, and $b^{-1}$ denotes the inverse of $b$ modulo $p$.

\subsection{The classes}
For an integer $a$ prime to $p$ let $\fq(a)\in\F_p$ be the Fermat quotient, $\fq(a)\equiv(a^{p-1}-1)/p\pmod p$.
It depends only on $a\bmod p^2$, and
\begin{equation}\label{eq:fq}
\fq(ab)=\fq(a)+\fq(b),\qquad \fq(1+pw)=-w,\qquad \fq(p-b)=\fq(b)+b^{-1}\quad(1\le b\le p-1);
\end{equation}
the last identity follows from $(p-b)^{p-1}\equiv b^{p-1}+p\,b^{p-2}\pmod{p^2}$.
For $c\in\F_p$ the function $\psi_c(a)=\zeta_p^{c\fq(a)}$ is a character of $U_{p^2}$, of order $p$ if $c\ne0$, and even since $\fq(-1)=0$.
Every character modulo $p^2$ is uniquely $\theta\psi_c$ with $\theta$ a character modulo $p$ (viewed modulo $p^2$), and $(\theta\psi_c)|_H=\psi_c|_H$ with $\psi_c(1+pw)=\zeta_p^{-cw}$.
Hence the restriction classes are
\[
C_c=\{\theta\psi_c:\ \theta\ \text{odd modulo }p\},\qquad c\in\F_p .
\]
The class $C_0$ consists of the odd characters modulo $p$, and $L(1,\theta)$ is the same for $\theta$ viewed modulo $p$ or modulo $p^2$; these values are linearly independent over $\Q$ by \cite{BBW73}.
By Theorem A, Baker's question for $p^2$ is therefore equivalent to the $\Q$-linear independence of each class $C_c$, $c\ne0$.

From now on $c\in\{1,\dots,p-1\}$.
For $1\le b\le p-1$ put
\[
z_b^{(c)}=\frac{\zeta_p^{-c\fq(b)}}{\zeta_p^{cb^{-1}}-1}\in\Q(\zeta_p),\qquad
e_c(\theta)=\bar\theta(c)\sum_{b=1}^{p-1}\theta(b)\,z_b^{(c)}\quad(\theta\text{ odd}).
\]

\begin{lemma}\label{lem:classformula}
For $\theta$ odd modulo $p$ and $c\in\{1,\dots,p-1\}$,
\[
\tau(\theta\psi_c)=p\,\theta(c)\,\zeta_p^{c\fq(c)}\zeta_{p^2}^{c},\qquad
B_{1,\bar\theta\psi_{-c}}=\sum_{b=1}^{p-1}\bar\theta(b)\,z_b^{(c)},\qquad
L(1,\theta\psi_c)=\frac{\pi i}{p}\,\zeta_p^{c\fq(c)}\zeta_{p^2}^{c}\;e_c(\bar\theta).
\]
In particular $C_c$ is linearly independent over $\Q$ if and only if the numbers $e_c(\theta)$, $\theta$ odd, are.
\end{lemma}

\begin{proof}
Let $\chi=\theta\psi_c$; it is primitive of conductor $p^2$.
Every $x\in U_{p^2}$ is uniquely $x=a(1+pw)$ with $1\le a\le p-1$ and $w\in\Z/p\Z$, and $\chi(x)e(x/p^2)=\chi(a)\zeta_p^{-cw}e(a/p^2)\zeta_p^{aw}$, where $e(t)=e^{2\pi it}$.
Summing over $w$ leaves only $a=c$, which gives $\tau(\chi)=p\chi(c)\zeta_{p^2}^c$ and the first formula.
For the second write $a=b+pk$ with $1\le b\le p-1$, $0\le k\le p-1$; then $a\equiv b(1+pkb^{-1})\pmod{p^2}$, so $\bar\chi(a)=\bar\theta(b)\zeta_p^{-c\fq(b)}\zeta_p^{ckb^{-1}}$, and $\sum_{k=0}^{p-1}k\,\eta^k=p/(\eta-1)$ for $\eta^p=1\ne\eta$ gives
\[
B_{1,\bar\chi}=\frac1{p^2}\sum_{b=1}^{p-1}\bar\theta(b)\zeta_p^{-c\fq(b)}\sum_{k=0}^{p-1}\zeta_p^{ckb^{-1}}(b+pk)=\sum_{b=1}^{p-1}\bar\theta(b)\,z_b^{(c)} .
\]
The third formula follows from \eqref{eq:Lprim}, since $\theta(c)\sum_b\bar\theta(b)z_b^{(c)}=e_c(\bar\theta)$.
The factor in front of $e_c(\bar\theta)$ does not depend on $\theta$, and $\theta\mapsto\bar\theta$ permutes the odd characters.
\end{proof}

\begin{lemma}\label{lem:symm}
$z^{(c)}_{p-b}=-z^{(c)}_b$ for $1\le b\le p-1$.
Consequently $e_c(\theta)=2\bar\theta(c)\sum_{b=1}^{(p-1)/2}\theta(b)z_b^{(c)}$ for odd $\theta$.
\end{lemma}

\begin{proof}
By \eqref{eq:fq}, $\fq(p-b)=\fq(b)+b^{-1}$ and $(p-b)^{-1}\equiv-b^{-1}$, so with $x=\zeta_p^{cb^{-1}}$ we get $z_{p-b}^{(c)}=\zeta_p^{-c\fq(b)}x^{-1}/(x^{-1}-1)=\zeta_p^{-c\fq(b)}/(1-x)=-z_b^{(c)}$.
\end{proof}

Lemma~\ref{lem:symm} is where the Fermat quotient of the \emph{integer} $p-b$, rather than of $-b$, matters: the twist by $\fq$ is not a function on $\F_p^\times$.

\subsection{A rational criterion}

\begin{theorem}\label{thm:star}
For $1\le b\le\frac{p-1}2$ let
\[
X_b=\frac{\zeta_p^{\fq(b)+b^{-1}}}{\zeta_p^{b^{-1}}-1}=\frac1p\sum_{j=0}^{p-1}\bigl\langle bj-b\fq(b)-1\bigr\rangle\,\zeta_p^{\,j},
\]
where $\langle x\rangle\in\{0,\dots,p-1\}$ is the least non-negative residue.
If the numbers $X_b$, $1\le b\le\frac{p-1}2$, are linearly independent over $\Q$, then Baker's question has an affirmative answer for $q=p^2$; moreover, the values in each class $C_c$, $c\ne0$, are then linearly independent over $\Q(\zeta_{p-1})$.
\end{theorem}

\begin{proof}
The expansion of $X_b$ follows from $1/(\zeta_p^{a}-1)=\frac1p\sum_{i=0}^{p-1}i\,\zeta_p^{ai}$ with $a=b^{-1}$, by substituting $j=a i+\fq(b)+b^{-1}$.
Let $c\ne0$ and let $x_\theta\in\Q(\zeta_{p-1})$, $\theta$ odd, satisfy $\sum_\theta x_\theta e_c(\theta)=0$.
The function $h=\sum_\theta x_\theta\theta$ on $\F_p^\times$ is odd, and $\sum_\theta x_\theta e_c(\theta)=\sum_{b=1}^{p-1}h(bc^{-1})z_b^{(c)}=2\sum_{b\le(p-1)/2}h(bc^{-1})z_b^{(c)}$ by Lemma~\ref{lem:symm}.
Let $(\alpha_k)$ be a $\Q$-basis of $\Q(\zeta_{p-1})$ and write $h(bc^{-1})=\sum_kh_k(b)\alpha_k$ with $h_k(b)\in\Q$.
Since $\Q(\zeta_p)$ and $\Q(\zeta_{p-1})$ are linearly disjoint over $\Q$ and $z_b^{(c)}\in\Q(\zeta_p)$, we get $\sum_{b\le(p-1)/2}h_k(b)z_b^{(c)}=0$ for every $k$.
Let $\sigma_s\in\Gal(\Q(\zeta_p)/\Q)$ be given by $\zeta_p\mapsto\zeta_p^s$.
Then $\sigma_{-c}(X_b)=-\sigma_{-c}\bigl(\zeta_p^{\fq(b)}/(\zeta_p^{-b^{-1}}-1)\bigr)=-z_b^{(c)}$, so $\sum_bh_k(b)X_b=0$ and hence $h_k(b)=0$ for all $k$ and all $b\le(p-1)/2$.
Thus $h$ vanishes on $\{bc^{-1}:1\le b\le\frac{p-1}2\}$, which together with its negatives is $\F_p^\times$; as $h$ is odd, $h=0$, and so all $x_\theta=0$.
The first assertion now follows from Theorem A.
\end{proof}

In coordinates with respect to the basis $1,\zeta_p,\dots,\zeta_p^{p-2}$, the criterion asks that the $\frac{p-1}2\times(p-1)$ integer matrix with rows $\bigl(\langle bj-b\fq(b)-1\rangle-\langle-b\fq(b)-1-b\rangle\bigr)_{0\le j\le p-2}$ have full rank; here $j=p-1$ was eliminated by $\zeta_p^{p-1}=-\sum_{j<p-1}\zeta_p^j$.
Without the Fermat quotient the $X_b$ would be $\frac12-\frac i2\cot(\pi b^{-1}/p)$, and the criterion would reduce to Chowla's theorem on cotangents.

\subsection{Galois descent}
The group $P=\{t\in U_{p(p-1)}:t\equiv1\pmod p\}\cong U_{p-1}$ of Lemma~\ref{lem:Psym} acts on the odd characters by $\theta\mapsto\theta^t$.
Concretely, for $t\in U_{p-1}$ let $\tau_t$ be the automorphism of $\Q(\zeta_{p(p-1)})$ with $\tau_t(\zeta_p)=\zeta_p$ and $\tau_t(\zeta_{p-1})=\zeta_{p-1}^t$; then $\tau_t(e_c(\theta))=e_c(\theta^t)$.
Let $p-1=2^vm$ with $m$ odd and
\[
D=\{d\mid p-1:\ 2^v\mid d\}.
\]
An odd character has order in $D$, and for $d\in D$ the odd characters of order $d$ form one orbit $\{\theta_d^u:u\in U_d\}$, where $\theta_d$ is any one of them and $\theta_d^u$ only depends on $u\bmod d$.
Let $M$ be the exponent of $U_{p-1}$ and $K_1=\Q(\zeta_{p-1},\zeta_M)$; all characters $\phi$ of $U_{p-1}$ take values in $K_1$, and $p$ is unramified in $K_1$.
For $d\in D$ let $G_d=\ker(U_{p-1}\to U_d)$, and for a character $\phi$ of $U_{p-1}$ let $D_\phi=\{d\in D:\ \phi|_{G_d}=1\}$.
For $d\in D_\phi$ the character $\phi$ factors through $U_d$, and we put
\[
\varepsilon^{\phi}_d(c)=\sum_{u\in U_d}\bar\phi(u)\,e_c(\theta_d^u).
\]

\begin{proposition}\label{prop:descent}
Let $c\ne0$. If for every character $\phi$ of $U_{p-1}$ the numbers $\varepsilon_d^\phi(c)$, $d\in D_\phi$, are linearly independent over $K_1$, then the class $C_c$ is linearly independent over $\Q$.
\end{proposition}

\begin{proof}
Let $O$ be the set of odd characters modulo $p$, and let $R\subseteq\Q^O$ be the space of vectors $(x_\theta)$ with $\sum_\theta x_\theta e_c(\theta)=0$.
Since $\tau_t$ fixes $\Q$, $R$ is stable under the permutation action $(t\cdot x)_{\theta^t}=x_\theta$ of $U_{p-1}$, and so is $R_1=R\otimes_\Q K_1\subseteq K_1^O$; every $x\in R_1$ still satisfies $\sum_\theta x_\theta e_c(\theta)=0$.
Suppose $R\ne0$.
Since $K_1$ contains the values of all characters of the abelian group $U_{p-1}$, the module $R_1$ is the sum of its $\phi$-isotypic parts, so there are $\phi$ and $0\ne x\in R_1$ with $t\cdot x=\phi(t)x$ for all $t$.
The permutation module $K_1^O$ is the direct sum over $d\in D$ of the modules $K_1^{O_d}$ of the orbits $O_d\cong U_{p-1}/G_d$, and the $\phi$-eigenspace of $K_1^{O_d}$ is zero unless $d\in D_\phi$, in which case it is spanned by $\sum_{u\in U_d}\bar\phi(u)[\theta_d^u]$.
Hence $x=\sum_{d\in D_\phi}y_d\sum_{u\in U_d}\bar\phi(u)[\theta_d^u]$ with $y_d\in K_1$ not all $0$, and $0=\sum_\theta x_\theta e_c(\theta)=\sum_{d\in D_\phi}y_d\,\varepsilon_d^\phi(c)$, contrary to the hypothesis.
\end{proof}

\subsection{\texorpdfstring{$\fP$}{P}-adic expansion}
Fix a prime $\fP$ of $L_1=K_1(\zeta_p)$ above $p$, let $\omega$ be the Teichm\"uller character, $\omega(b)\equiv b\pmod\fP$, and let $\varpi=\zeta_p-1$.
Let $\nu$ be the valuation at $\fP$ normalised by $\nu(p)=1$.
Since $p$ is unramified in $K_1$ and $L_1/K_1$ is totally ramified at $\fP$ of degree $p-1$, we have $\nu(\varpi)=1/(p-1)$ and $\nu(K_1^\times)=\Z$.
Let $\Ok$ be the valuation ring of the completion of $L_1$ at $\fP$; then $\Ok/p\Ok=\kappa[\varpi]/(\varpi^{p-1})$ with $\kappa$ the residue field, so an element of $\Ok$ that is non-zero modulo $p$ has valuation $i/(p-1)$, $i$ the index of its first non-zero coefficient in $\kappa[\varpi]/(\varpi^{p-1})$.
For odd $j$ with $1\le j\le p-4$ put
\[
\beta_j=\frac{B_{j+1}}{j+1}\bmod p,\qquad \gamma_j=\frac12\sum_{b=1}^{p-1}b^{j+1}\fq(b)^2\bmod p .
\]

\begin{lemma}\label{lem:sums}
Let $k$ be even with $2\le k\le p-3$.
Then $\sum_{b=1}^{p-1}b^k\fq(b)\equiv-B_k/k\pmod p$, and $B_{1,\omega^{k-1}}\equiv B_k/k\pmod\fP$.
In particular $\beta_j\equiv-\sum_bb^{j+1}\fq(b)\equiv B_{1,\omega^j}\pmod\fP$.
\end{lemma}

\begin{proof}
Let $S_n=\sum_{b=1}^{p-1}b^n$.
For even $n\ge2$ with $(p-1)\nmid n$, Faulhaber's formula
\[
S_n=\frac1{n+1}\sum_{i=0}^{n}\binom{n+1}{i}B_ip^{n+1-i}=\sum_{i=0}^{n}\binom ni\frac{p^{n+1-i}}{n+1-i}B_i
\]
gives $S_n\equiv pB_n\pmod{p^2}$: the term $i=n$ is $pB_n$ with $B_n$ $p$-integral, the term $i=n-1$ is $\frac n2p^2B_{n-1}$, and for $i\le n-2$ the term with $r=n+1-i\ge3$ has valuation at least $r-v_p(r)-1\ge2$ (as $v_p(B_i)\ge-1$ and $p\ge5$).
Hence $\sum_bb^k\fq(b)=(S_{k+p-1}-S_k)/p\equiv B_{k+p-1}-B_k\pmod p$, and Kummer's congruence $B_{k+p-1}/(k+p-1)\equiv B_k/k$ gives $B_{k+p-1}\equiv\frac{k-1}{k}B_k$, whence the first claim.
For the second, $a^p$ is a $(p-1)$-st root of unity modulo $p^2$ congruent to $a$ modulo $p$, so $\omega(a)\equiv a^p\pmod{p^2}$, and with $n=p(k-1)+1$ we get $pB_{1,\omega^{k-1}}=\sum_a\omega^{k-1}(a)a\equiv S_n\equiv pB_n\pmod{p^2}$.
As $n\equiv k\pmod{p-1}$ and $n\equiv1\pmod p$, Kummer's congruence gives $B_n\equiv nB_k/k\equiv B_k/k\pmod p$.
\end{proof}

\begin{lemma}\label{lem:expansion}
Let $c\in\{1,\dots,p-1\}$.
\begin{enumerate}
\item $\varpi\,e_c(\omega^{-1})\equiv-1\pmod\fP$; hence $\nu(e_c(\omega^{-1}))=-1/(p-1)$.
\item For odd $j$ with $1\le j\le p-4$, $e_c(\omega^j)$ is $\fP$-integral and $e_c(\omega^j)\equiv c^{-j}\beta_j\pmod\fP$.
\item For odd $j$ with $3\le j\le p-4$,
\[
\varpi\,e_c(\omega^j)\equiv\omega(c)^{-j}\bigl(\tilde\beta_j\,\varpi+c\,\tilde\gamma_j\,\varpi^2\bigr)\pmod{\varpi^3\Ok},
\]
where $\tilde\beta_j,\tilde\gamma_j\in\Z$ are any lifts of $\beta_j,\gamma_j$.
\end{enumerate}
\end{lemma}

\begin{proof}
Fix $b$, let $k\in\{1,\dots,p-1\}$ with $k\equiv cb^{-1}$ and let $n\ge0$ be an integer with $n\equiv-c\fq(b)\pmod p$.
In $\Z_p[[X]]$ we have $(1+X)^n=1+nX+\binom n2X^2+\cdots$ and
\[
\frac{X}{(1+X)^k-1}=\frac1{k+\binom k2X+\binom k3X^2+\cdots}=\frac1k\Bigl(1-\frac{k-1}2X+\frac{k^2-1}{12}X^2+\cdots\Bigr),
\]
both series having coefficients in $\Z_p$ since $p\ge5$ and $p\nmid k$.
Evaluating at $X=\varpi$, so that $(1+\varpi)^n=\zeta_p^{-c\fq(b)}$ and $(1+\varpi)^k=\zeta_p^{cb^{-1}}$, we get
\begin{equation}\label{eq:Pb}
\varpi z_b^{(c)}=\frac1k\Bigl(1+\bigl(n-\tfrac{k-1}2\bigr)\varpi+\bigl(\tbinom n2-\tfrac{n(k-1)}2+\tfrac{k^2-1}{12}\bigr)\varpi^2\Bigr)+O(\varpi^3),
\end{equation}
with $O(\varpi^3)\subseteq\varpi^3\Ok$.
Since $p\in\varpi^{p-1}\Ok\subseteq\varpi^3\Ok$, each coefficient may be replaced by anything congruent to it modulo $p\Ok$; in particular $\omega^j(b)$ by $b^j$ (indeed $\omega(b)-b$ lies in $\fP$ and in the unramified ring $\Z_p[\mu_{p-1}]$, hence in $p\Z_p[\mu_{p-1}]$), $1/k$ by $b/c$, $k$ by $c/b$ and $n$ by $-c\fq(b)$.
Multiplying \eqref{eq:Pb} by $\omega^j(b)$ and summing over $b$, the coefficient of $\varpi^0$ is $\sum_b\omega^j(b)/k\equiv c^{-1}\sum_bb^{j+1}\pmod p$, which is $\equiv-c^{-1}$ if $j=p-2$ and $\equiv0$ otherwise.
For $j=p-2$ this gives (1), because $\bar\omega^{p-2}(c)=\omega(c)\equiv c$.
For $j\le p-4$ the constant coefficient lies in $p\Ok\subseteq\varpi^3\Ok$; this gives the integrality in (2).
The coefficient of $\varpi$ is congruent to
\[
\frac1c\sum_bb^{j+1}\Bigl(-c\fq(b)-\frac{c/b-1}2\Bigr)=-\sum_bb^{j+1}\fq(b)-\frac12\sum_bb^j+\frac1{2c}\sum_bb^{j+1}\equiv\beta_j
\]
by Lemma~\ref{lem:sums}, since $(p-1)\nmid j$ and $(p-1)\nmid j+1$; with $\bar\omega^j(c)\equiv c^{-j}$ this proves (2).
The coefficient of $\varpi^2$ is congruent to
\[
\frac1c\sum_bb^{j+1}\Bigl(\frac{c^2\fq(b)^2}2+\frac{c^2\fq(b)}{2b}+\frac{c^2}{12b^2}-\frac1{12}\Bigr)
=c\gamma_j+\frac c2\sum_bb^{j}\fq(b)+\frac c{12}\sum_bb^{j-1}-\frac1{12c}\sum_bb^{j+1}.
\]
For $3\le j\le p-4$ the last two sums vanish modulo $p$, and pairing $b$ with $p-b$ and using \eqref{eq:fq} gives $2\sum_bb^j\fq(b)\equiv-\sum_bb^{j-1}\equiv0$.
This proves (3).
\end{proof}

Consequently, for every $d\in D$ with $d<p-1$ the numbers $\varepsilon_d^\phi(c)$ are $\fP$-integral, while for the orbit of $\omega^{-1}$, which has order $p-1\in D$, taking $\theta_{p-1}=\omega^{-1}$ we get from (1) and (2)
\begin{equation}\label{eq:polar}
\nu\bigl(\varepsilon_{p-1}^\phi(c)\bigr)=-\frac1{p-1}\qquad\text{for every }\phi\text{ and every }c\ne0 .
\end{equation}

\subsection{Proof of Theorem C}
We first record a parity property of Bernoulli numbers of quartic characters.

\begin{lemma}\label{lem:parity}
Let $p\equiv5\pmod 8$, $p\ge13$, and let $\theta$ be a character of order $4$ modulo $p$ (it is odd).
Then $B_{1,\theta}=x+iy$ with odd integers $x,y$, where $i=\theta(g)$ for a primitive root $g$.
\end{lemma}

\begin{proof}
We have $pB_{1,\theta}=\sum_a\theta(a)a\in\Z[i]$, and $B_{1,\theta}$ is integral at both primes of $\Q(i)$ above $p$: Lemma~\ref{lem:sums} applies with $\fP$ above either of them, since $\theta$, of order $4<p-1$, is not $\omega^{-1}$ for either choice. Hence $x,y\in\Z$.
Let $F$ be the imaginary cyclic quartic field of conductor $p$.
By the analytic class number formula \cite[Theorem~4.17]{Was97}, $h^-(F)=Qw\prod_{\chi\in\{\theta,\bar\theta\}}\bigl(-\tfrac12B_{1,\chi}\bigr)$, where $w=2$ is the number of roots of unity in $F$ and $Q=1$ is the unit index, so $x^2+y^2=|B_{1,\theta}|^2=2h^-(F)$.
(To see $Q=1$: for a unit $\eta$ of $F$, $\eta/\bar\eta$ is a root of unity, i.e.\ $\pm1$, and $\bar\eta=-\eta$ is impossible, since then $F=F^+(\sqrt{-\eta\bar\eta})$ with $\eta\bar\eta$ a unit of the real quadratic subfield $F^+$, and $F/F^+$ would be unramified at $p$.)
The class number of $F$ is odd: in the cyclic extension $F/\Q$ of degree $4$ the ramified places are $p$ (with $e=4$) and $\infty$ (with $e=2$), and $-1$ is not a norm from the totally imaginary field $F$, so Chevalley's ambiguous class number formula \cite{Che33} gives $|\mathrm{Cl}(F)^{\Gal(F/\Q)}|=\frac{4\cdot2}{4\cdot2}=1$; a $2$-group acting on a non-trivial finite $2$-group has a non-trivial fixed point, so the $2$-part of $\mathrm{Cl}(F)$ is trivial.
Since $h^-(F)$ divides $h(F)$ \cite[Theorem~4.10]{Was97}, $h^-(F)$ is odd, and $x^2+y^2\equiv2\pmod4$ forces $x$ and $y$ to be odd.
\end{proof}

For $p-1=2^v\ell$ with $\ell$ an odd prime and $v\ge2$ we define the sums, for $\xi\in\mu_{2^v}(\F_p)$ and homomorphisms $\phi:U_{2^v}\to\F_p^\times$,
\[
S_\beta(\xi,\phi)=\sum_{u\in U_{2^v}}\phi(u)^{-1}\xi^u\beta_{\ell u},\qquad
S_\gamma(\xi,\phi)=\sum_{u\in U_{2^v}}\phi(u)^{-1}\xi^u\gamma_{\ell u},
\]
where $u$ runs over the odd integers in $[1,2^v)$ (so that $3\le\ell u\le p-1-\ell$), and the condition
\[
(\mathrm C_p)\qquad\bigl(S_\beta(\xi,\phi),S_\gamma(\xi,\phi)\bigr)\ne(0,0)\quad\text{for all }\xi\in\mu_{2^v}(\F_p)\text{ and all }\phi .
\]

\begin{theorem}\label{thm:psquare}
Let $p\ge5$ be a prime.
Baker's question has an affirmative answer for $q=p^2$ in each of the following cases:
\begin{enumerate}
\item $p$ is a Fermat prime;
\item $p=2\ell+1$ with $\ell$ an odd prime;
\item $p=4\ell+1$ with $\ell$ an odd prime;
\item $p=2^v\ell+1$ with $v\ge3$ and $\ell$ an odd prime, and $(\mathrm C_p)$ holds.
\end{enumerate}
\end{theorem}

\begin{proof}
By Lemma~\ref{lem:classformula} and Proposition~\ref{prop:descent} it suffices to show, for every $c\ne0$ and every character $\phi$ of $U_{p-1}$, that the numbers $\varepsilon_d^\phi(c)$, $d\in D_\phi$, are linearly independent over $K_1$.
Note $p-1\in D_\phi$ for every $\phi$.
If $D_\phi=\{p-1\}$, this holds by \eqref{eq:polar}.
If $D_\phi=\{d,p-1\}$ with $d<p-1$, it suffices to show that $\varepsilon_d^\phi(c)\ne0$ and $\nu(\varepsilon_d^\phi(c))\notin-\frac1{p-1}+\Z$: a relation $y_1\varepsilon_d^\phi(c)+y_2\varepsilon_{p-1}^\phi(c)=0$ with $y_1,y_2\in K_1$ not both zero would force both to be non-zero and $\nu(\varepsilon^\phi_d(c))-\nu(\varepsilon^\phi_{p-1}(c))=\nu(y_2/y_1)\in\Z$.

(1) Here $p-1=2^v$, so $D=\{p-1\}$.

(2) Here $p\equiv3\pmod 4$, $D=\{2,p-1\}$ and $G_2=U_{p-1}$, so $D_\phi=\{2,p-1\}$ for $\phi=1$ and $D_\phi=\{p-1\}$ otherwise.
For $\phi=1$ we have $\varepsilon_2^1(c)=e_c(\rho)$ with $\rho=\omega^{\ell}$ the quadratic character, and by Lemmas~\ref{lem:expansion}(2) and~\ref{lem:sums}, $e_c(\rho)\equiv c^{-\ell}\beta_\ell\equiv\rho(c)B_{1,\rho}=-\rho(c)h(-p)\pmod\fP$, where $h(-p)$ is the class number of $\Q(\sqrt{-p})$.
Since $0<h(-p)=|B_{1,\rho}|<\frac1p\sum_{a=1}^{p-1}a<p$, $e_c(\rho)$ is a unit at $\fP$, and $\nu(e_c(\rho))=0$.

(3) and (4)
Here $D=\{2^v,p-1\}$ and $G_{2^v}=\{t\equiv1\pmod{2^v}\}$, so $D_\phi=\{2^v,p-1\}$ if $\phi$ factors through $U_{2^v}$ and $D_\phi=\{p-1\}$ otherwise.
Take $\theta_{2^v}=\omega^{\ell}$; for $u\in U_{2^v}$ represented by an odd integer $1\le u<2^v$, $\theta_{2^v}^u=\omega^{\ell u}$ with $3\le\ell u\le p-1-\ell\le p-4$.
By Lemma~\ref{lem:expansion}(3), with $\xi=c^{-\ell}\in\mu_{2^v}(\F_p)$ and $\phi$ reduced modulo $\fP$ (as $c$ runs over $\F_p^\times$, $\xi$ runs over all of $\mu_{2^v}(\F_p)$, and the reductions of the characters of $U_{2^v}$ are all the homomorphisms $U_{2^v}\to\F_p^\times$, so $(\mathrm C_p)$ is exactly what is needed),
\[
\varpi\,\varepsilon_{2^v}^\phi(c)\equiv S_\beta(\xi,\phi)\,\varpi+c\,S_\gamma(\xi,\phi)\,\varpi^2\pmod{(p,\varpi^3)}
\]
(coefficients read in $\kappa\supseteq\F_p$).
Under $(\mathrm C_p)$, $\nu(\varpi\varepsilon_{2^v}^\phi(c))\in\{\frac1{p-1},\frac2{p-1}\}$, so $\nu(\varepsilon^\phi_{2^v}(c))\in\{0,\frac1{p-1}\}$, and neither value lies in $-\frac1{p-1}+\Z$ because $p\ge5$.
This proves (4), and it proves (3) once we show $S_\beta(\xi,\phi)\ne0$ when $v=2$.
For $v=2$, $U_4=\{1,3\}$ and $S_\beta(\xi,\phi)=\xi\bigl(\beta_\ell+\phi(3)^{-1}\xi^2\beta_{3\ell}\bigr)$ with $\phi(3)^{-1}\xi^2=\pm1$.
By Lemma~\ref{lem:sums}, $\beta_\ell\equiv B_{1,\theta}$ and $\beta_{3\ell}\equiv B_{1,\theta^3}=B_{1,\bar\theta}$ modulo $\fP$, where $\theta=\omega^\ell$ has order $4$.
With $B_{1,\theta}=x+iy$ as in Lemma~\ref{lem:parity} we get $\beta_\ell+\beta_{3\ell}\equiv2x$ and $\beta_\ell-\beta_{3\ell}\equiv2iy$, and these are non-zero modulo $\fP$ because $x,y$ are odd and $|x|,|y|\le|B_{1,\theta}|<\frac{p-1}2$.
\end{proof}

\begin{remark}[On condition $(\mathrm C_p)$]\label{rem:Cp}
The first sums $S_\beta$ alone do not suffice when $v\ge3$.
For $p=41$ one finds $B_{1,\theta}=-1+i-\sqrt2$ for a suitable character $\theta$ of order $8$ (the other three are its Galois conjugates), whose component along $\sqrt{-2}$ vanishes; correspondingly $S_\beta(\xi,\phi)=0$ for eight pairs $(\xi,\phi)$ with $\phi$ of conductor $8$.
Such vanishing is an archimedean accident: the coordinates of $B_{1,\theta}$ are small integers.
The second coefficients $\gamma_j$ are of a different nature; by the Kummer congruences modulo $p^2$ one has $\gamma_j\equiv-\frac1p\bigl(\frac{B_{j+p}}{j+p}-\frac{B_{j+1}}{j+1}\bigr)\pmod p$, a discrete derivative of $B_n/n$ along $n\equiv j+1\pmod{p-1}$.
We have checked $(\mathrm C_p)$ for all primes $p<1.6\cdot10^5$ of the form $2^v\ell+1$ with $v\ge3$ (Section~\ref{sec:comp}); there are 624 of them, $(\mathrm C_p)$ holds for all of them, and the first sums alone fail for 180 of them.
\end{remark}

\begin{remark}[Relation with Iwasawa power series]\label{rem:iwasawa}
By Lemma~\ref{lem:classformula}, $e_c(\theta)$ is, up to the factor $\bar\theta(c)$, the Bernoulli number $B_{1,\theta\psi_{-c}}$, and $\psi_{-c}(1+p)=\zeta_p^{c}$.
By \cite[Theorem~7.10]{Was97}, $B_{1,\theta\psi}=-f(\zeta_\psi-1,\theta\omega)$ for characters $\psi$ of the second kind, where $f(T,\theta\omega)\in\Z_p[\mu_{p-1}][[T]]$ is the Iwasawa power series and, in Washington's normalisation, $\zeta_\psi=\psi(1+p)^{-1}$; thus $\zeta_{\psi_{-c}}=\zeta_p^{-c}$; for $\theta=\omega^{-1}$, i.e.\ $\theta\omega=1$, the formula carries in addition the denominator of the $p$-adic zeta function, which vanishes at $T=p$ and has valuation $1/(p-1)$ at $T=\zeta_p^{-c}-1$; this is the polar term of Lemma~\ref{lem:expansion}(1).
Thus, up to the factors $-\bar\theta(c)$, the class $C_c$ consists of the values at $T=\zeta_p^{-c}-1$ of the power series attached to the even characters $\theta\omega$, and the coefficients of $\varpi,\varpi^2,\dots$ in Lemma~\ref{lem:expansion} are governed by the first, second, \dots\ Iwasawa coefficients modulo $p$.
Theorem~\ref{thm:psquare} uses the constant terms (Bernoulli numbers), the class numbers $h(-p)$ and $h^-(F)$, and, in case (4), the linear terms.
For general $p$ the descent produces sets $D_\phi$ with up to $\tau(m)$ elements ($p-1=2^vm$), and separating them requires either more coefficients or more refined information on the orbits. There are general results on Iwasawa power series modulo $p$, such as the vanishing of the $\mu$-invariant \cite{FW79} and the linear independence modulo $p$ of the Iwasawa series of characters of conductor prime to $p$ \cite{AR10}, but they do not control the particular values and coefficients needed here, and we do not see how to remove this obstacle.
\end{remark}
\section{Computations}\label{sec:comp}

All computations use exact integer arithmetic (Python with \texttt{numpy} and \texttt{python-flint}); the programs and their output are in the supplementary archive.

\subsection{Certificates for restriction classes}

\begin{lemma}\label{lem:certificate}
Let $r\equiv1\pmod N$ be a prime, $g\in\F_r^\times$ of order $N$, and for $j\in U_N$ let $\rho_j:\Z[\zeta_N,1/N]\to\F_r$ be the ring homomorphism with $\rho_j(\zeta_N)=g^j$.
Let $C$ be a restriction class and $J\subseteq U_N$.
If the matrix $\bigl(\rho_j(W_\chi)\bigr)_{j\in J,\ \chi\in C}$ has rank $|C|$ over $\F_r$, then the numbers $L(1,\chi)$, $\chi\in C$, are linearly independent over $\Q$.
If $j\equiv j'\pmod\lambda$, the rows of $j$ and $j'$ are proportional.
\end{lemma}

\begin{proof}
We have $W_\chi\in q^{-1}\Z[\zeta_N]\subseteq\Z[\zeta_N,1/N]$.
If $\sum_\chi x_\chi W_\chi=0$ with coprime integers $x_\chi$, then $\sum_\chi x_\chi\rho_j(W_\chi)=0$ for all $j\in J$, so the reduction of $(x_\chi)$ modulo $r$ lies in the kernel of the matrix and is $0$, a contradiction.
If $j'=ju$ with $u\equiv1\pmod\lambda$, then $\rho_{j'}=\rho_j\circ\sigma_u$ with $\sigma_u\in\Gamma$, and $\sigma_u$ multiplies every $W_\chi$, $\chi\in C_\mu$, by the same root of unity $\bar\mu(u\bmod q)$ (Lemma~\ref{lem:Gamma}).
\end{proof}

By Theorem A it suffices to certify each restriction class separately.
For each of the $2708$ moduli $3\le q\le10^4$ with $\varphi(\delta)\ge4$ we took for $r$ the least prime $r\equiv1\pmod N$ above $2^{24}$ (above $2^{40}$ if $N\ge2^{23}$), one $j$ in each residue class of $U_N$ modulo $\lambda$ (in random order), and added rows until every class had full rank.
All $964{,}083$ restriction classes of these moduli (all of size $\varphi(\delta)/2\ge2$, the largest of size $729$ for $q=3^8$) reached full rank; the computation took about $35$ minutes on one core.
Together with Corollary B this proves the first part of Theorem D.

As an independent check, which does not use Theorem A, we computed for every $3\le q\le200$ the rank over $\Z$ of the vectors $qW_\chi\in\Z[\zeta_N]\cong\Z^{\varphi(N)}$ ($\chi$ odd), written in the power basis after reduction modulo the cyclotomic polynomial $\Phi_N$; in all cases it equals the number of odd characters.

\subsection{Prime squares}
For every prime $p<5000$ we computed the rank modulo the prime $2^{61}-1$ of the $\frac{p-1}2\times(p-1)$ integer matrix of Theorem~\ref{thm:star}; it is always $\frac{p-1}2$.
Full rank modulo a prime implies full rank over $\Q$, so Theorem~\ref{thm:star} gives the second part of Theorem D.

For condition $(\mathrm C_p)$ we computed $\beta_{\ell u}$ and $\gamma_{\ell u}$ from their definitions as sums of Fermat quotients (Lemma~\ref{lem:sums}), and the tables of all sums $S_\beta(\xi,\phi)$ and $S_\gamma(\xi,\phi)$ by exact matrix products.
Table~\ref{tab:Cp} summarises the result for all primes $p<1.6\cdot10^5$ of the form $2^v\ell+1$ with $\ell$ an odd prime and $v\ge2$.

\begin{table}[ht]
\centering
\begin{tabular}{crrr}
\toprule
$v$ & primes & all $S_\beta\ne0$ & $(\mathrm C_p)$ holds\\
\midrule
2 & 514 & 514 & 514\\
3 & 277 & 258 & 277\\
4 & 150 & 120 & 150\\
5 & 90 & 52 & 90\\
6 & 44 & 13 & 44\\
7 & 25 & 1 & 25\\
8 & 18 & 0 & 18\\
9 & 9 & 0 & 9\\
10 & 4 & 0 & 4\\
11 & 3 & 0 & 3\\
12 & 2 & 0 & 2\\
13 & 1 & 0 & 1\\
14 & 1 & 0 & 1\\
\midrule
all & 1138 & 958 & 1138\\
\bottomrule
\end{tabular}
\caption{Condition $(\mathrm C_p)$ for the primes $p=2^v\ell+1<1.6\cdot10^5$, $\ell$ an odd prime.}\label{tab:Cp}
\end{table}

For $v=2$ the first sums never vanish, as predicted by Lemma~\ref{lem:parity}.
We also validated the ingredients: $\beta_j$ agrees with $B_{j+1}/(j+1)\bmod p$ computed from Bernoulli numbers for all $p<400$; for all $p<1200$ we computed $\varpi e_1(\omega^j)$ in $\F_p[\varpi]/(\varpi^{24})$ directly from the definition and confirmed the constant and linear coefficients in Lemma~\ref{lem:expansion} for all odd $j$ (and the quadratic coefficients for $p<300$ and several $c$), and for the primes $p<1200$ of the form $2^v\ell+1$, $v\ge1$, the valuations of all $\varepsilon^\phi_{2^v}(c)$ obtained in this way are $0$ or $\frac1{p-1}$, in agreement with the computation of $(\mathrm C_p)$; and for the $568$ primes $p\equiv5\pmod 8$ below $2\cdot10^4$ both coordinates of $B_{1,\theta}$, $\theta$ quartic, are odd, as Lemma~\ref{lem:parity} asserts.
\section{Remarks on the remaining cases}\label{sec:remarks}

By Theorem A, a $\Q$-linear relation among the odd values can only live inside a restriction class $C_\mu=\chi_0E_\delta$, and these classes are non-trivial exactly when $\varphi(\delta)\ge4$.
A prime $p$ divides $\delta$ either because $p^2\mid q$ (we call the corresponding part of the class \emph{wild}), or because $p\parallel q$ and $p\mid r-1$ for some prime $r\mid q$ (\emph{tame}).
Section~\ref{sec:psquare} treats the simplest wild case.
In the simplest tame case Stickelberger's theorem separates part of the classes.

\begin{proposition}\label{prop:tame}
Let $q=pq'$ with $p$ an odd prime not dividing $q'$, and suppose that $\delta=p$.
Then every restriction class has the form $C=\{\alpha\chi':\ \alpha\bmod p,\ \alpha(-1)=-\chi'(-1)\}$ for a fixed character $\chi'$ modulo $q'$.
If the order of $\chi'$ is prime to $p$, the values $L(1,\chi)$, $\chi\in C$, are linearly independent over $\Q$.
\end{proposition}

\begin{proof}
Here $H=\{a\equiv1\pmod p\}$ and $E_\delta$ is the group of even characters modulo $p$; writing $\chi=\alpha\chi'$ with $\alpha$ modulo $p$ and $\chi'$ modulo $q'$, the restriction $\chi|_H$ determines $\chi'$, which gives the form of $C$.
Let $n$ be the order of $\chi'$, $\chi'^*$ the primitive character of conductor $f'\mid q'$ inducing $\chi'$, $F_0=\Q(\mu_{p-1},\mu_n,\zeta_{q'})$ and $F=F_0(\zeta_p)$.
Since $p\nmid nq'(p-1)$, the field $F_0$ is unramified at $p$; let $\fP$ be a prime of $F$ above $p$, $\nu$ the valuation with $\nu(p)=1$ (so $\nu(F_0^\times)=\Z$), and $\omega$ the Teichm\"uller character at $\fP$.
If $\alpha=\omega^{-a}$ with $1\le a\le p-2$, then $\chi=\alpha\chi'$ is induced by $\alpha\chi'^*$, of conductor $pf'$, with $\tau(\alpha\chi'^*)=\alpha(f')\chi'^*(p)\tau(\alpha)\tau(\chi'^*)$, and by \eqref{eq:Lprim}
\[
\frac{L(1,\chi)}{\pi i}=\frac{\alpha(f')\chi'^*(p)\tau(\alpha)\tau(\chi'^*)}{pf'}\,B_{1,\bar\alpha\bar\chi'^*}\prod_{r\mid q',\ r\nmid f'}\Bigl(1-\frac{\alpha(r)\chi'^*(r)}r\Bigr).
\]
All factors except $\tau(\alpha)$ are non-zero elements of $F_0$, and by Stickelberger's congruence \cite[\S6.2]{Was97} $\nu(\tau(\omega^{-a}))=a/(p-1)$.
Hence $\nu(L(1,\chi)/\pi i)\equiv a/(p-1)\pmod\Z$.
If $\alpha=1$ (which happens when $\chi'$ is odd), then $L(1,\chi)=L(1,\chi'^*)(1-\chi'^*(p)/p)\prod_{r}(1-\chi'^*(r)/r)\in\pi i\,F_0^\times$ and $\nu(L(1,\chi)/\pi i)\in\Z$.
So the numbers $L(1,\chi)/\pi i$, $\chi\in C$, have pairwise different valuations modulo $\Z$, and in a non-trivial relation $\sum_\chi x_\chi L(1,\chi)/\pi i=0$ with $x_\chi\in\Q$ the term of smallest valuation could not cancel.
\end{proof}

When $p$ divides the order of $\chi'$, the values of $\chi'$ contain $p$-th roots of unity, the Bernoulli factors lie in a field ramified at $p$, and the valuations of the class elements are no longer determined modulo $\Z$ by the Gauss sums.
For instance, if $q=pr$ with $r\equiv1\pmod p$ prime and $\chi'$ has $p$-power order, then the reduction of $\chi'$ modulo $\fP$ is trivial, and the corresponding ``leading term'' $B_{1,\bar\alpha}$ acquires the Euler factor $1-\bar\alpha(r)=0$; the classes are then governed by finer invariants, reminiscent of the trivial zeros in Gross's refined conjecture on Stickelberger elements \cite{Gro88}.
For the wild classes of higher prime powers $p^k$, $k\ge3$, the classes have $(p-1)p^{k-2}/2$ elements and involve characters of order $p^{k-1}$; the expansion of Section~\ref{sec:psquare} then takes place in powers of $\zeta_{p^{k-1}}-1$ and needs correspondingly more coefficients.
We have not been able to treat these cases in general; in the range $q\le10^4$ they are covered by the computations of Section~\ref{sec:comp}.
\begin{thebibliography}{99}

\bibitem{AR10}
B.~Angl\`es and G.~Ranieri, On the linear independence of $p$-adic $L$-functions modulo $p$, \emph{Ann. Inst. Fourier (Grenoble)} \textbf{60} (2010), 1831--1855.

\bibitem{Bak75}
A.~Baker, \emph{Transcendental Number Theory}, Cambridge University Press, Cambridge, 1975.

\bibitem{BBW73}
A.~Baker, B.~J.~Birch and E.~A.~Wirsing, On a problem of Chowla, \emph{J. Number Theory} \textbf{5} (1973), 224--236.

\bibitem{Che33}
C.~Chevalley, Sur la th\'eorie du corps de classes dans les corps finis et les corps locaux, \emph{J. Fac. Sci. Univ. Tokyo Sect.~I} \textbf{2} (1933), 365--476.

\bibitem{Cho70}
S.~Chowla, The nonexistence of nontrivial linear relations between the roots of a certain irreducible equation, \emph{J. Number Theory} \textbf{2} (1970), 120--123.

\bibitem{FW79}
B.~Ferrero and L.~C.~Washington, The Iwasawa invariant $\mu_p$ vanishes for abelian number fields, \emph{Ann. of Math.} \textbf{109} (1979), 377--395.

\bibitem{GK22}
S.~Gun and N.~Kandhil, On an extension of a question of Baker, preprint, arXiv:2212.00376 (2022).

\bibitem{Gro88}
B.~H.~Gross, On the values of abelian $L$-functions at $s=0$, \emph{J. Fac. Sci. Univ. Tokyo Sect.~IA Math.} \textbf{35} (1988), 177--197.

\bibitem{KR22}
N.~Kandhil and P.~Rath, Around a question of Baker, preprint, arXiv:2212.00358 (2022).

\bibitem{MM11}
M.~Ram Murty and V.~Kumar Murty, A problem of Chowla revisited, \emph{J. Number Theory} \textbf{131} (2011), 1723--1733.

\bibitem{Oka80}
T.~Okada, On an extension of a theorem of S.~Chowla, \emph{Acta Arith.} \textbf{38} (1980/81), 341--345.

\bibitem{Was97}
L.~C.~Washington, \emph{Introduction to Cyclotomic Fields}, 2nd ed., Graduate Texts in Mathematics~83, Springer, New York, 1997.

\end{thebibliography}
\end{document}
